\section{مدل \lr{Hodgkin--Huxley}}
\subsection*{تولید \lr{action potential}}
این مدل درنظر می‌گیرد که در داخل نورون پتاسیم بیشتر است و در خارج نورون سدیم. وقتی یک محرک باعث می‌شود تا نورون دیپلاریزه شود، و ولتاژ آن از حالت استراحت (حدود $-70$ میلی ولت) به حد آستانه (حدود $-55$ میلی ولت) برسد گیت‌های سدیم شروع به باز شدن می‌کنند. این مورد باعث می‌شود از آنجایی که غلظت سدیم در بیرون غشا بیشتر است، یون‌های سدیم به داخل سلول بیایند و نورون بیشتر دیپلاریزه شود.  سپس نورون به حد بیشینه ولتاژ خود (حدود $+35$ میلی ولت) می‌رسد و در اینجا دیگر گیت‌های سدیم بسته می‌شوند و گیت‌های پتاسیم باز می‌شوند تا دوباره ولتاژ نورون به حالت استراحت برگردد. بعد از اینکه نورون به حالت استراحت برمی‌گردد این گیت‌ها با کمی تاخیر بسته می‌شوند در نتیجه نورون کمی هایپر پلاریزه می‌شود و سپس به حالت استراحت برمی‌گردد. 
\subsection*{نقش پتاسیم}
پتاسیم در این مدل این وظیفه را دارد که پس از اسپایک زدن نورون، ولتاژ نورون را به حالت استراحت بازگرداند. همچنین چون بعد از اسپایک زدن گیت‌های پتاسیم باز هستند، پس اسپایک در جهت رو به عقب نمی‌تواند حرکت کند زیرا قسمت عقب وارد فاز 
\lr{refactory period}
شده و یون‌های پتاسیم باعث می‌شوند تا پتانسیل آن قسمت از حد آستانه بیشتر نشود پس اسپایک رو به عقب نمی‌آید.
\subsection*{\lr{refactory period} در مدل \lr{HH}}
این موضوع به اینصورت نشان داده می‌شود که پس از اسپایک زدن نورون، گیت‌های سدیم بسته می‌شوند و گیت‌های پتاسیم باز می‌شوند. و نورون به سمت پلاریزه شدن می‌رود.  درواقع فاز از بعد از اسپایک زدن تا رسیدن ولتاژ نورون به حالت استراحت 
\lr{refactory period}
است.
\subsection*{محدودیت‌های مدل \lr{HH}}
این مدل علیرغم دستاوردش محدودیت‌های زیادی دارد. از جمله:
\begin{itemize}
	\item 
	این مدل تنها بر اساس یک نورون خاص است. سایر نورون‌ها در بقیه بدن لزوما به اینصورت نیستند و پارامترهای آن‌ها ممکن است فرق کند.
	\item 
	این مدل فقط گیت‌های سدیم و پتاسیم را درنظر می‌گیرد و از نقش سایر یون‌ها مثل کلسیم و کلر صرف نظر می‌کند. 
	\item 
	این مدل فقط تاثیر گیت‌ها را درنظر می‌گیرد و درباره سیناپس، نوروترنسمیتور و همچنین تقویت سیناپسی و یادگیری  اطلاعاتی نمی‌دهد.
	\item 
	از لحاظ محاسباتی باید سه معادله دیفرانسیل را حل کرد و سپس جواب آن ها را در یک معادله جبری گذاشت تا اختلاف پتانسیل نورون بدست بیاید. یعنی از لحاظ محاسباتی پیچیده است. 
\end{itemize}
\subsection*{نقش کانال‌های یونی در مدل}
این مدل اختلاف پتانسیل نورون را ناشی از برآیند سه نوع کانال یونی می‌بیند. پس اولا باید دید این کانال‌ها با چه احتمالی باز هستند  (که خود ناشی از ریت باز و بسته شدن گیت ها می‌باشد) و ثانیا اختلاف غلظت بین داخل سلول و بیرون آن چقدر است. هرچقدر گیت‌های بیشتری باز باشند و اختلاف غلظت بیشتری موجود باشد، روند تغییر پتانسیل نیز بیشتر است. 
\subsection*{بهبودهای روی مدل \lr{HH}}
برای حل مشکلاتی که ذکر شد، یکی از راه‌ها ساده‌سازی مدل بوده که با ثابت درنظر گرفتن احتمال باز بودن گیت سدیم (به علت سرعت زیاد آن) و تخمین زدن گیت‌های 
\lr{refactory}
سدیم برابر با یک منهای گیت‌های پتاسیم، دو مدل از دستگاه کم کرد. همچنین یک سری مدل‌ها به دنبال این بوده اند که به مدل 
\lr{HH}
یون‌ها و کانال‌های جدید اضافه کنند تا نورون‌های بیشتری را بتوان با آن تحلیل کرد.
\subsection{آزمایش‌های شبیه‌سازی}
































	

	
	\subsection{تولید پتانسیل عمل}
	
	\subsubsection{تعیین جریان آستانه}
	برای تعیین جریان آستانه مورد نیاز برای ایجاد پتانسیل عمل، از روش جستجوی دودویی با پالس جریان به مدت ۱ میلی‌ثانیه استفاده شد. ن
	
	\textbf{جریان آستانه:} تقریباً \textbf{$6.84$ میکروآمپر بر سانتی‌متر مربع} برای پالس ۱ میلی‌ثانیه‌ای
	
	\begin{figure}[H]
		\centering
		\includegraphics[width=0.8\textwidth]{pics/1_threshold_staircase.png}
		\caption{
			حداکثر پتانسیل غشا بر حسب دامنه جریان اعمال‌شده. }
		\label{fig:threshold_staircase}
	\end{figure}
	همانطور که در شکل 
	\ref{fig:threshold_staircase}
	می‌توان دید قبل از حد آستانه ولتاژ نورون تقریبا رفتار غیر خطی دارد. هنگام رسیدن به آستانه یک دفعه پرش شدیدی دارد و بعد از حد آستانه تقریبا ثابت است.
	
	\begin{figure}[H]
		\centering
		\includegraphics[width=0.8\textwidth]{pics/2_threshold_traces.png}
		\caption{مقدار پتانسیل غشا بر حسب زمان در جریان‌های متفاوت حول آستانه.}
		\label{fig:threshold_traces}
	\end{figure}
	
	\subsubsection{اثرات دامنه و مدت جریان}
	پتانسل غشا و به طور خاص مدت زمان اسپایک زدن تحت تاثیر دامنه و مدت زمان جریان است.
	
	\textbf{اثرات دامنه:}
	\begin{itemize}
		\item افزایش دامنه جریان، زمان تا شروع پیک را کاهش می‌دهد
		\item  دامنه‌های بالاتر باعث می‌شوند تا نورون با شیب سریعتری اسپایک بزند اما دامنه پیک تقریبا مشابهی دارند.
	\end{itemize}
	
	\begin{figure}[H]
		\centering
		\includegraphics[width=0.8\textwidth]{pics/3_amplitude_effect.png}
		\caption{شکل‌های موج پتانسیل عمل برای دامنه‌های جریان مختلف (پالس ۱ میلی‌ثانیه‌ای). .}
		\label{fig:amplitude_effect}
	\end{figure}
	
	\textbf{اثرات مدت:}
	به نظر می‌آید مدت نیز اثری شبیه به دامنه دارد. یعنی باعث می‌شود شیب اکشن پتانسیل بیشتر شود. همچنین با دادن جریان بصورت دائم نورون بصورت متناوب اسپایک می‌زند. 

	
	\begin{figure}[H]
		\centering
		\includegraphics[width=0.8\textwidth]{pics/4_duration_effect.png}
		\caption{
			پتانسیل نورون هنگام اعمال جریان ۱۰ میکرو آمپر برای مدت‌های مختلف.
			}
		\label{fig:duration_effect}
	\end{figure}
	
	\textbf{تحلیل فرکانس اسپایک زدن:}
	همانطور که در شکل 
	\ref{fig:fi_curve}
	می‌توان دید
	 با افزایش جریان فرکانس اسپایک زدن تقریبا بصورت خطی زیاد می‌شود. 
	
	\begin{figure}[H]
		\centering
		\includegraphics[width=0.8\textwidth]{pics/5_fi_curve.png}
		\caption{رابطه بین فرکانس اسپایک و جریان }
		\label{fig:fi_curve}
	\end{figure}
	
	\subsubsection{اثر شرایط اولیه}
در این قسمت شرایط اولیه مختلف برای نورون بررسی می‌شود. 

این شرایط شامل ولتاژ اولیه مختلف و نرخ باز بودن گیت‌های مختلف است. 
	
\begin{table}[H]
	\centering
	\caption{مقایسه پتانسیل برای شرایط اولیه مختلف}
	\label{tab:initial_conditions}
	\begin{latin}
\begin{tabular}{lcccccc}
	\hline
	\textbf{ Condition} & \textbf{Threshold (mV)} & \textbf{Peak (mV)} & \textbf{Afterhyperpolarization (mV)}  & \textbf{Peak Time (ms)} \\
	\hline
	Default              & -50 & 39 & -76 & 2.6 \\
	Depolarized          & -47 & 35 & -76  & 2.2 \\
	Hyperpolarized       & -53 & 41 & -76  & 2.9 \\
	High $n$             & - & - & - & -  \\
	High $m$             & -46 & 40 & -77 & 1.8 \\
	Low $h$              & - & - & - & - \\
	\hline
\end{tabular}
	\end{latin}
\end{table}

	\begin{figure}[H]
		\centering
		\includegraphics[width=\textwidth]{pics/6_initial_conditions.png}
		\caption{اختلاف پتانسیل نورون به ازای شرایط اولیه مختلف}
		\label{fig:initial_conditions}
	\end{figure}
	
	\textbf{یافته‌های کلیدی:}
	\begin{itemize}
		\item شرایط اولیه دپولاریزه زمان رسیدن به پیک را کاهش می‌دهد
		\item شرایط هیپرپولاریزه زمان نهفتگی را افزایش و دامنه را کمی افزایش می‌دهد
		\item شکل پتانسیل عمل در برابر تغییرات متوسط در شرایط اولیه مقاوم است
		\item \lr{h} پایین یا \lr{n} بالا باعث جلوگیری از ایجاد پیک می‌شوند.
	\end{itemize}
	
	\subsection{رسانایی کانال و دینامیک}
	
	\subsubsection{اثر رسانایی سدیم }
	مقدار رسانایی سدیم دامنه و شروع پتانسیل عمل را تعیین می‌کند. (درواقع منظور رسانایی بیشینه است وگرنه در هر لحظه رسانایی بر اساس معادلات دیفرانسیل تعیین می‌شود.)
	
	\begin{figure}[H]
		\centering
		\includegraphics[width=0.8\textwidth]{pics/7_gNa_effect.png}
		\caption{
			اثر تغییر $g_{Na}$ بر شکل موج پتانسیل عمل.
			}
		\label{fig:gna_effect}
	\end{figure}
	
	\textbf{مشاهدات:}
	\begin{itemize}
		\item افزایش  $g_{Na}$ دامنه پیک و شیب فاز صعودی را افزایش می‌دهد
		\item انتظار می‌رفت با کاهش رسانایی نورون اسپایک نزند ولی در آزمایش‌های من حتی با مقدار ۱۰ و جریان کمتر از حد یافت شده در بخش اول همچنان نورون اسپایک می‌زند.
		\item عرض پیک با افزایش  $g_{Na}$کاهش می‌یابد
	\end{itemize}
	
	\subsubsection{اثر رسانایی پتاسیم ( $g_{K}$)}

	
	\begin{figure}[H]
		\centering
		\includegraphics[width=0.8\textwidth]{pics/8_gK_effect.png}
		\caption{
			اثر تغییر $g_{K}$ بر شکل ولتاژ نورون.
			}
		\label{fig:gk_effect}
	\end{figure}
	
	\textbf{مشاهدات:}
	\begin{itemize}
		\item افزایش $g_{K}$  مدت پتانسیل عمل را کوتاه می‌کند
		\item $g_{K}$  پایین منجر به پیک‌های پهن‌تر و دوره‌های مقاوم طولانی‌تر می‌شود
		\item $g_{K}$  پایین باعث می‌شود تا نورون چندبار اسپایک بزند. 
	\end{itemize}
	
	\subsubsection{اثرات ترکیبی $g_{Na}$ و $g_{K}$}
در این بخش تاثیر دو بخش قبل بصورت ترکیبی نمایش داده می‌شود.
	
	\begin{figure}[H]
		\centering
		\includegraphics[width=0.8\textwidth]{pics/9_gNa_gK_map.png}
		\caption{
			مقدار فرکانس اسپایک زدن بر حسب ترکیبات $g_{Na}$  و $g_{K}$ با جریان ۱۰ میکروآمپر بر سانتی‌متر مربع. رنگ مقدار فرکانس را نشان می‌دهد.
			}
		\label{fig:gna_gk_map}
	\end{figure}
	
	\textbf{یافته‌های کلیدی:}
	\begin{itemize}
		\item $g_{Na}$ بالا و $g_{K}$پایین به فرکانس بالای اسپایک زدن کمک می‌کنند
		\item افزایش هر دو رسانایی می‌تواند فرکانس‌های اسپایک زدن مشابهی را حفظ کند
		\item نسبت $g_{Na}$/$g_{K}$ آستانه تحریک‌پذیری را تعیین می‌کند
		\item ترکیبات مختلف می‌توانند پتانسیل‌های عمل کیفی مشابهی تولید کنند
	\end{itemize}
	
	\subsubsection{نقش رسانایی نشتی ($g_{L}$)}
	رسانایی نشتی، پتانسیل استراحت غشا را تعیین کرده و بر مقاومت ورودی تأثیر می‌گذارد.
	
	\begin{figure}[H]
		\centering
		\includegraphics[width=\textwidth]{pics/10_gL_effect.png}
		\caption{
			اثر $g_{L}$ بر پتانسیل استراحت (چپ) و پتانسیل عمل (راست). 
			}
		\label{fig:gl_effect}
	\end{figure}
	
	\textbf{مشاهدات:}
	\begin{itemize}
		\item افزایش  پتانسیل استراحت غشا را دپولاریزه می‌کند
		\item  هرچه $g_L$ بزرگ تر باشد نورون راحت تر اسپایک می‌زند. 
		\item  $g_L$ خیلی پایین باعث می‌شود نورون پیک نزند.
	\end{itemize}
	این موضوع با انتظار من در تضاد است. زیرا انتظار داشتم 
	$g_L$ 
	باعث شود تا جریان تخلیه شود و نورون اسپایک نزند. با جستجو به نظر می‌آید این مورد به این خاطر است که 
	$g_L$ 
	پایین باعث می‌شود تا گیت‌های یونی فعال و یا غیرفعال شوند و به همین جهت روی اسپایک زدن نورون تاثیر می‌گذارد.
	\subsection{حساسیت پارامترها}
	
	\subsubsection{تغییرات سیستماتیک پارامترها}
	تحلیل حساسیت جامعی با تغییر هر پارامتر و ثابت نگه‌داشتن سایر پارامترها در مقادیر پیش‌فرض انجام شد.
	
	\begin{figure}[H]
		\centering
		\includegraphics[width=\textwidth]{pics/11_sensitivity.png}
		\caption{تحلیل حساسیت نشان‌دهنده اثر تغییر $g_{Na}$، ، $g_{L}$ و ظرفیت غشا  $C$ بر ولتاژ پیک، عرض پتانسیل عمل و تعداد پیک‌ها.}
		\label{fig:sensitivity}
	\end{figure}
	

	
	\subsubsection{پیامدها برای دینامیک عصبی}
	تحلیل حساسیت نشان می‌دهد که مدل هاچکین-هاکسلی به ویژه به پارامترهای کنترل‌کننده رسانایی‌های سدیم و پتاسیم حساس است که با نقش بنیادی آنها در تولید پتانسیل عمل همسو است:
	
	\begin{enumerate}
		\item \textbf{
			کنترل تحریک‌پذیری:} $g_{Na}$ تعیین‌کننده اصلی تحریک‌پذیری است و می‌توان دید یون اصلی در اسپایک زدن نورون همین یودن سدیم است. 
		

		\item \textbf{ظرفیت نورون و $g_L$}
		این موارد تاثیر کمتری در اکشن پتانسیل نورون دارند. ظرفیت نورون که عمدتا یک مقدار ثابت است و در ازمایش‌ها $g_L$  نسبت به سدیم و پتاسیم خیلی اثر کمتری دارد.
	\end{enumerate}
	
	\subsection{رفتار زیرآستانه}
	
	\subsubsection{پاسخ‌های پتانسیل غشای زیرآستانه}
در شکل 
\ref{fig:subthreshold_traces}
به نورون جریان‌های کمتر از حد آستانه‌ای که در بخش اول یافت شد داده شده و ولتاژ آن رسم شده است.
	
	\begin{figure}[H]
		\centering
		\includegraphics[width=0.8\textwidth]{pics/12_subthreshold_traces.png}
		\caption{پاسخ‌های پتانسیل غشای زیرآستانه به جریان‌های پایدار. .}
		\label{fig:subthreshold_traces}
	\end{figure}
	همانطور که در شکل
	\ref{fig:subthreshold_traces}
	می‌توان دید هنگامی که جریان باعث اسپایک زدن نورون نمی‌شود، نورون وارد یک فاز 
	\lr{refactory}
	می‌شود و سپس به حالت ثابت می‌رسد. (علیرغم این که نمودار شبیه به نمودار اسپایک زدن است ولی با دقت به مقیاس ولتاژ می‌توان دید که پتانسیل نورون تنها حدود دو میلی ولت فرق می‌کند. )
	
	\subsubsection{رابطه ورودی-خروجی}
	رابطه بین دامنه جریان زیرآستانه و دپولاریزاسیون حالت پایدار رفتار تقریباً خطی را نشان می‌دهد.
	
	\begin{figure}[H]
		\centering
		\includegraphics[width=0.8\textwidth]{pics/13_subthreshold_depol.png}
		\caption{دپولاریزاسیون حالت پایدار بر حسب دامنه جریان زیرآستانه. .}
		\label{fig:subthreshold_depol}
	\end{figure}
	
	\textbf{مشاهدات کلیدی:}
	\begin{itemize}
		\item رابطه خطی برای جریان‌های کوچک
			\item با نزدیک شدن به حد آستانه رابطه از خطی بودن خارج می‌شود. 
		\item مقاومت ورودی را می‌توان از شیب محاسبه کرد
		\item اشباع نشان‌دهنده رسانایی‌های وابسته به ولتاژ است ، مثلا گیت‌های \lr{volume gated}
	\end{itemize}
	
	\subsubsection{نقش متغیرهای \lr{n, m, h}  در رفتار زیرآستانه}
	متغیرهای 
	\lr{n, m, h}
	 دینامیک‌های متمایزی را در طول تحریک زیرآستانه نشان می‌دهند.
	
	\begin{figure}[H]
		\centering
		\includegraphics[width=0.8\textwidth]{pics/14_subthreshold_gating.png}
		\caption{
			پتانسیل غشای زیرآستانه و متغیرهای 
			\lr{n, m, h}
			 در طول تحریک جریان زیرآستانه پایدار.}
		\label{fig:subthreshold_gating}
	\end{figure}
	همانطور که در شکل 
	\ref{fig:subthreshold_gating}
	می‌توان دید، پس از اعمال جریان بیرونی، گیت‌های سدیم باز می‌شوند تا ولتاژ را به حالت تعادل بازگرداند. این اختلاف پتانسیل جزئی باعث می‌شود تا گیت‌های پتاسیم نیز مقداری باز بشوند و درنهایت نورون پایدار می‌شود. 
	
	هرچقدر مقدار این جریان بیشتر باشد گیت‌های سدیم و پتاسیم بیشتری باز می‌شوند پس یعنی 
	\lr{n}
	و
	\lr{m}
	زیاد می‌شوند و 
	\lr{h}
	کم می‌شود .
